\section{Related work} \label{sec:related}

We review graph neural networks for fraud detection, explainability methods for GNNs and how they are evaluated, temporal graph learning, and low-latency deployment. We then position RTXGNN against the closest methods.

\subsection{Graph neural networks for fraud detection}

Graph Neural Networks model the relational structure of transaction networks. Graph Convolutional Networks (GCN) \citep{kipf2017gcn, asiri2025graph}, Graph Attention Networks (GAT) \citep{velickovic2018gat,fu2026phishing} and GraphSAGE \citep{hamilton2017graphsage,tang2025application} established the message-passing framework on which most fraud detectors build. GCN aggregates neighbours with a normalised adjacency, GAT learns attention weights over neighbours, and GraphSAGE learns inductive aggregators that generalise to unseen nodes, which matters when new accounts or transactions arrive continuously.

Fraud graphs violate several assumptions of these architectures. \textbf{Camouflage}, in which fraudsters mimic legitimate users, weakens neighbourhood aggregation. CARE-GNN \citep{dou2020caregnn} selects neighbours with a label-aware similarity measure and a reinforcement-learning threshold. \textbf{Class imbalance} is addressed by PC-GNN \citep{liu2021pcgnn} with label-balanced sampling of training nodes and neighbour selection. FRAUDRE \citep{zhang2021fraudre} tackles both graph inconsistency and imbalance with a fraud-aware convolution that uses node--neighbour differences, a combination of intermediate layer outputs, and an imbalance-oriented classifier. GAS \citep{li2019gas}, deployed at Xianyu, combines aggregation over the interaction graph with aggregation over a similarity graph built from content features. \citet{liu2020graphconsis} characterise context, feature and relation inconsistencies and propose GraphConsis. H2-FDetector \citep{shi2022h2fdetector} separates homophilic from heterophilic edges, GeniePath \citep{liu2019geniepath} learns adaptive receptive fields, and DGA-GNN \citep{duan2024dgagnn} groups non-additive attributes with decision-tree binning. \citet{tang2022rethinking} relate fraud and anomaly detection on graphs to spectral "right-shift" effects and release the T-Finance dataset used in our study. GADBench \citep{tang2023gadbench} shows that tree ensembles with neighbourhood aggregation are strong competitors to specialised GNNs on many anomaly-detection benchmarks.

On the Elliptic Bitcoin data specifically, \citet{weber2019anti} found that a random forest on the transaction features outperforms GCN, and that all models fail after a dark-market shutdown in the test period. This is consistent with the broader observation that tree ensembles remain very strong on tabular features \citep{grinsztajn2022tree}. We therefore include logistic regression, random forest \citep{breiman2001rf} and XGBoost \citep{chen2016xgboost} as baselines, in addition to general, temporal and fraud-specific GNNs.

\subsection{Explainability methods for graph neural networks}

\textbf{Post-hoc explainers} explain a trained model. GNNExplainer \citep{ying2019gnnexplainer} optimises soft masks over edges and features that preserve a prediction and must be run separately for every instance. PGExplainer \citep{luo2020pgexplainer} amortises this optimisation with a parametric edge-mask network. SubgraphX \citep{yuan2021subgraphx} searches subgraphs with Monte Carlo tree search and Shapley values and is accurate but slow. CF-GNNExplainer \citep{lucic2022cfgnnexplainer} finds minimal edge deletions that change a prediction, and CF$^2$ \citep{tan2022counterfactual} combines factual (sufficient) and counterfactual (necessary) criteria. GraphLIME \citep{huang2022graphlime} fits local surrogate models, XGNN \citep{yuan2020xgnn} gives model-level explanations, and gradient-based attributions such as saliency \citep{simonyan2014saliency} and Integrated Gradients \citep{sundararajan2017ig} apply to GNNs directly \citep{pope2019explainability, baldassarre2019explainability}. \citet{yuan2023explainability_survey} survey the field.

\textbf{Self-explaining GNNs} compute explanations during the forward pass. GSAT \citep{miao2022gsat} learns stochastic attention under an information bottleneck. SEFraud \citep{li2024sefraud} learns feature and edge masks for fraud detection on heterogeneous graphs with a triplet loss that separates the masked from the complementary representation, and reports deployment at the Industrial and Commercial Bank of China. SEFraud and GSAT operate on static graphs. RTXGNN adopts the self-explaining design but uses an objective that ties the reported top-$k$ explanation to the prediction through explicit sufficiency and necessity terms, and it adds relative temporal encodings. We use the term "self-explaining" rather than "interpretable": a network with several non-linear layers remains a complex model even when it exposes attributions \citep{rudin2019stop, lipton2018mythos}.

\textbf{Evaluating explanations} is difficult. \citet{faber2021groundtruth} show that comparing explanations with synthetic ground truth can mislead, because a model may solve a task with a rationale different from the planted one. Fidelity metrics measure instead how the prediction changes when the explanation is removed (Fid$^+$, necessity) or kept alone (Fid$^-$, sufficiency) \citep{yuan2021subgraphx, amara2022graphframex}. \citet{agarwal2023graphxai} add stability and fairness criteria. We report Fid$^+$ and Fid$^-$ at several explanation sizes, with a random explanation as reference. We also report stability under input perturbation and agreement across independently trained models, class-wise results, and the cost of each explainer. Ground truth is used only on the synthetic benchmark, where the planted rings are known to be the only signal separating the classes.

\subsection{Temporal graph neural networks}

Discrete-time models operate on snapshots. EvolveGCN \citep{pareja2020evolvegcn} evolves the GCN weights with a recurrent network (EvolveGCN-O/-H), DySAT \citep{sankar2020dysat} applies structural and temporal self-attention, and ROLAND \citep{you2022roland} turns static GNNs into incremental dynamic models. Continuous-time models consume timestamped events. TGAT \citep{xu2020tgat} encodes time gaps with a functional time encoding inside temporal attention. TGN \citep{rossi2020tgn} adds a per-node memory updated by a recurrent cell from event messages and shows that JODIE, DyRep and TGAT are special cases. APAN \citep{wang2021apan} decouples inference from graph queries with asynchronous mailboxes, a design aimed at real-time serving. Time2Vec \citep{kazemi2019time2vec} is a generic learnable time embedding with one linear and several periodic components.

Continuous-time models assume that nodes interact repeatedly over time. In Elliptic each transaction exists in one of 49 snapshots and all edges connect transactions of the same snapshot \citep{weber2019anti}. The memory of TGN and the mailbox of APAN therefore carry no information across snapshots, and time gaps within a snapshot are zero. We still include TGAT, TGN, APAN and EvolveGCN as baselines, as requested by the reviewers, and describe their adaptation in \ref{app:baselines}. DySAT's temporal attention runs over the snapshots in which a node appears. Every Elliptic transaction appears in exactly one snapshot, so DySAT reduces to a structural attention network (GAT), which is already a baseline.

\subsection{Low-latency fraud detection systems}

Production fraud detection systems have end-to-end latency budgets of roughly 10--100\,ms. BRIGHT \citep{lu2022bright}, deployed at eBay, pre-computes entity embeddings in a batch layer and runs only a light real-time layer for new transactions. TitAnt \citep{cao2019titant,dong2025enhancing} and GEM \citep{liu2018gem, rao2025fraud} serve predictions at Ant Financial and Alipay. These systems return scores without explanations. When explanations are needed, they are usually produced asynchronously for the cases that are escalated. Latency figures reported in the literature depend strongly on hardware, batch size and whether neighbourhood retrieval is included. We therefore report our own end-to-end measurements on commodity CPU hardware, with every stage of the pipeline timed separately.

Beyond graph data, supervised learning is widely used for credit-card and insurance fraud. \citet{afriyie2023supervised} compare logistic regression, decision trees and random forests for credit-card fraud and find the random forest strongest under class imbalance. \citet{jacob2025golden} combine optimisation-based feature selection with a BERT-LSTM model for insurance claims. These results reinforce the need to compare graph models against strong feature-based learners.

\subsection{Regulation and explanation requirements}

GDPR Article~22\footnote{\url{https://eur-lex.europa.eu/eli/reg/2016/679/oj}} restricts solely automated decisions with legal or similarly significant effects, and Articles~13--15 require "meaningful information about the logic involved". Whether this amounts to a right to an explanation is debated \citep{wachter2017righttoexplanation, selbst2017meaningful, goodman2016explanation, sargeant2026mind}. In U.S. consumer credit, the Equal Credit Opportunity Act \citep{ecoa1974} and Regulation~B require specific principal reasons for adverse actions. The CFPB has stated that this obligation also applies when complex algorithms are used\footnote{\url{https://www.consumerfinance.gov/compliance/circulars/circular-2022-03-adverse-action-notification-requirements-in-connection-with-credit-decisions-based-on-complex-algorithms/}}. For virtual assets, FATF guidance\footnote{\url{https://www.fatf-gafi.org/en/publications/Fatfrecommendations/Guidance-rba-virtual-assets-2021.html}} and the EU MiCA regulation \citep{teng2026digital} require AML/CFT controls that are documented and auditable. Supervisory guidance on model risk management, such as the U.S. SR~11-7 \citep{sr117}, requires that models be validated and their limitations understood.

These texts set requirements for processes and for the content of notices, not for model architectures. A feature attribution can contribute to a reason statement or an audit record, but it does not by itself establish compliance with any of them. This applies especially when the features are anonymised, as in Elliptic. Section~\ref{subsec:regulatory} discusses what would be needed to operationalise RTXGNN's explanations. \citet{cherif2024survey} note that GNN explanations have not yet been validated with analysts or regulators, and our work does not close this gap.

\subsection{Positioning of RTXGNN}

Table~\ref{tab:related_comparison} compares RTXGNN with the closest methods on properties that can be verified from the methods themselves. We have removed the "regulatory compliance" column of the first version, because compliance is a property of an institution's process and cannot be verified from a method. We also replaced "real-time explanation" with "explanation in the forward pass", since latency depends on the deployment.

\begin{table}[!t]
\centering
\caption{Properties of related methods. \textbf{F} = supported, \textit{P} = partially supported, $-$ = not supported. Temporal: uses timestamps or snapshots. Self-expl.: explanation computed in the forward pass. Fraud: design for camouflage or imbalance. Feat./Edge: explanation over input features / edges. Obj.: explanation tied to the prediction by a sufficiency/necessity objective.}
\label{tab:related_comparison}
\footnotesize
\resizebox{\textwidth}{!}{%
\begin{tabular}{lcccccc}
\toprule
\textbf{Method} & \textbf{Temporal} & \textbf{Self-expl.} & \textbf{Fraud} & \textbf{Feat.} & \textbf{Edge} & \textbf{Obj.}\\
\midrule
GCN / GraphSAGE \citep{kipf2017gcn,hamilton2017graphsage} & $-$ & $-$ & $-$ & $-$ & $-$ & $-$\\
GAT \citep{velickovic2018gat} & $-$ & \textit{P} & $-$ & $-$ & \textit{P} & $-$\\
CARE-GNN / PC-GNN \citep{dou2020caregnn,liu2021pcgnn} & $-$ & $-$ & \textbf{F} & $-$ & $-$ & $-$\\
FRAUDRE / GAS \citep{zhang2021fraudre,li2019gas} & $-$ & $-$ & \textbf{F} & $-$ & $-$ & $-$\\
EvolveGCN \citep{pareja2020evolvegcn} & \textbf{F} & $-$ & $-$ & $-$ & $-$ & $-$\\
TGAT / TGN / APAN \citep{xu2020tgat,rossi2020tgn,wang2021apan} & \textbf{F} & \textit{P} & $-$ & $-$ & \textit{P} & $-$\\
GNNExplainer \citep{ying2019gnnexplainer} (post hoc) & $-$ & $-$ & $-$ & \textbf{F} & \textbf{F} & \textit{P}\\
PGExplainer \citep{luo2020pgexplainer} (post hoc) & $-$ & $-$ & $-$ & $-$ & \textbf{F} & \textit{P}\\
GSAT \citep{miao2022gsat} & $-$ & \textbf{F} & $-$ & $-$ & \textbf{F} & \textit{P}\\
SEFraud \citep{li2024sefraud} & $-$ & \textbf{F} & \textbf{F} & \textbf{F} & \textbf{F} & \textit{P}\\
\midrule
\textbf{RTXGNN (ours)} & \textit{P} & \textbf{F} & \textit{P} & \textbf{F} & \textbf{F} & \textbf{F}\\
\bottomrule
\end{tabular}}
\par\smallskip\raggedright\footnotesize
RTXGNN's temporal support is marked partial because it uses relative time only, which is uninformative on Elliptic (Section~\ref{subsec:hrape}). Its fraud-specific support is partial because it relies on class-balanced training rather than on a camouflage-specific mechanism. Attention weights (GAT, TGAT) are marked partial because attention is not necessarily faithful to the prediction.
\end{table}
